\section{Conclusion}

% In this work, we introduce Colossal-AUTO to automate intra-op parallel training for large models. We build a solid and extensible infrastructure to provide needed information for automatic parallelization including \textbf{a)} a \textit{symbolic profiler} which could obtain memory and computing overhead in a trivial time, \textbf{b)} a \textit{cluster detector} which collects cluster hardware performance and topology, and \textbf{c)} a \textit{tensor layout manager} which uses a heuristic function to find an efficient conversion path in a reasonable time. Our system could find an optimal execution plan with our 2-stage solver and recompile it into a module instance. Finding an efficient distributed training strategy is a task that required extensive expert experience. We believe our system could lower the barrier for large model training for users without distributed training experience.

In this work, we introduce Colossal-Auto to jointly search for intra-op parallelism and activation checkpointing to generate the near-optimal distributed training plan for large-scale models. We use an analyzer to gather essential information about the hardware performance and the computation graph cost, employ a two-stage solver to find the near-optimal execution plan, and finally recompile the PyTorch model to a distributed model. 
Finding an efficient distributed training strategy is a challenging task that required extensive expert experience. 
We believe our system could lower the barrier for large model training for users without distributed training experience. 
Our future work could improve the robustness and versatility of this system further.